% DFFRNT AI Assistant — Installer Guide
% Overleaf-compatible (pdfLaTeX). Styling adapted from the DFFRNT Beta
% Prototype Report theme: Palatino body, garnet (uOttawa-inspired) accent,
% with markdown-style Note / Tip / Warning callouts for a user-guide feel.
\documentclass[11pt,letterpaper]{article}

%------------------------------------------------------------------ encoding --
\usepackage[T1]{fontenc}
\usepackage[utf8]{inputenc}

% Silence the harmless "Command \showhyphens has changed" notice microtype
% emits under recent LaTeX kernels.
\usepackage{silence}
\WarningFilter{latex}{Command}

%--------------------------------------------------------------------- fonts --
\usepackage{mathpazo}                 % Palatino body + matching math
\usepackage[scaled=0.92]{helvet}      % Helvetica-like sans (headings)
\usepackage{microtype}                % refined justification & spacing
\linespread{1.08}
\renewcommand{\sfdefault}{\rmdefault} % unify on Palatino, matching the report

%-------------------------------------------------------------------- layout --
\usepackage[letterpaper,margin=1in,headheight=22pt,headsep=18pt,footskip=28pt]{geometry}
\usepackage{tocloft}                  % load BEFORE parskip
\usepackage{parskip}
\usepackage{enumitem}
\setlist{itemsep=2pt,topsep=4pt,parsep=0pt}

%--------------------------------------------------------------------- colour --
\usepackage[table]{xcolor}
\definecolor{accent}{HTML}{6A24B5}    % DFFRNT brand purple
\definecolor{accentdk}{HTML}{4E1A86}  % deeper purple (headings / links)
\definecolor{brandlime}{HTML}{CFEF3C} % chartreuse accent
\definecolor{pagebg}{HTML}{E7E2F0}    % lavender title-page background
\definecolor{icnk}{HTML}{171634}      % dark navy icon tile
\definecolor{blocklt}{HTML}{D9CEEE}   % light lavender text on the purple block
\definecolor{blocklt2}{HTML}{C7B9E6}  % lighter lavender (block small print)
\definecolor{ink}{HTML}{1F2933}       % near-black body ink
\definecolor{inkmute}{HTML}{56616C}   % muted grey for metadata
\definecolor{rulegrey}{HTML}{C7CDD4}  % hairline rules
\definecolor{briefbg}{HTML}{F1ECF8}   % light lavender callout / code background

% Admonition palette (markdown-style callouts)
\definecolor{noteblue}{HTML}{2B6CB0}
\definecolor{notebg}{HTML}{EEF4FB}
\definecolor{tipgreen}{HTML}{2F7A55}
\definecolor{tipbg}{HTML}{EDF6F0}
\definecolor{warnred}{HTML}{B02A37}
\definecolor{warnbg}{HTML}{FBEEEF}

%------------------------------------------------------------- code listings ---
\usepackage{listings}
\lstset{
  basicstyle=\ttfamily\small,
  keywordstyle=\color{accentdk}\bfseries,
  commentstyle=\color{inkmute}\itshape,
  stringstyle=\color{accent},
  frame=single,
  rulecolor=\color{rulegrey}, framesep=6pt,
  showstringspaces=false, breaklines=true, tabsize=2,
  backgroundcolor=\color{briefbg!40},
  xleftmargin=4pt,
}

%---------------------------------------------------- section / TOC styling ---
\usepackage{titlesec}
\titleformat{\section}
  {\sffamily\LARGE\bfseries\color{accentdk}}
  {\makebox[1.6em][l]{\thesection}}{0em}{}
  [\vspace{3pt}{\color{accent}\titlerule[1pt]}]
\titleformat{\subsection}
  {\sffamily\Large\bfseries\color{ink}}
  {\makebox[2.0em][l]{\thesubsection}}{0em}{}
\titleformat{\subsubsection}
  {\sffamily\large\bfseries\color{ink}}
  {\thesubsubsection}{0.6em}{}
% FAQ questions: garnet, run-in
\titleformat{\paragraph}[runin]
  {\sffamily\bfseries\color{accentdk}}{}{0em}{}
\titlespacing*{\section}      {0pt}{0pt}{14pt}
\titlespacing*{\subsection}   {0pt}{16pt}{6pt}
\titlespacing*{\subsubsection}{0pt}{12pt}{4pt}
\titlespacing*{\paragraph}    {0pt}{8pt}{0.7em}

\renewcommand{\thesection}{\arabic{section}}
\renewcommand*\contentsname{Table of Contents}
\renewcommand{\cfttoctitlefont}{\sffamily\LARGE\bfseries\color{accentdk}}
\renewcommand{\cftsecfont}{\sffamily\bfseries\color{ink}}
\renewcommand{\cftsecpagefont}{\sffamily\bfseries\color{ink}}
\renewcommand{\cftsubsecfont}{\color{inkmute}}
\renewcommand{\cftsubsecpagefont}{\color{inkmute}}
\setlength{\cftbeforesecskip}{6pt}
\renewcommand{\cftsecleader}{\cftdotfill{\cftdotsep}}

%-------------------------------------------------------- header / footer -----
\usepackage{fancyhdr}
\pagestyle{fancy}
\fancyhf{}
\renewcommand{\headrulewidth}{0.4pt}
\renewcommand{\headrule}{\hbox to\headwidth{\color{rulegrey}\leaders\hrule height \headrulewidth\hfill}}
\renewcommand{\footrulewidth}{0pt}
\renewcommand{\sectionmark}[1]{\markboth{\thesection\quad #1}{}}
\fancyhead[L]{\footnotesize\sffamily\color{inkmute}Installer Guide}
\fancyhead[R]{\footnotesize\sffamily\color{inkmute}\nouppercase{\leftmark}}
\fancyfoot[R]{\footnotesize\sffamily\color{inkmute}Page \thepage}
\fancyfoot[L]{\footnotesize\sffamily\color{inkmute}DFFRNT Air-Gapped AI Assistant}
\fancypagestyle{plain}{%
  \fancyhf{}%
  \renewcommand{\headrulewidth}{0pt}%
  \fancyfoot[R]{\footnotesize\sffamily\color{inkmute}Page \thepage}%
  \fancyfoot[L]{\footnotesize\sffamily\color{inkmute}DFFRNT Air-Gapped AI Assistant}%
}

%------------------------------------------------------- hyperlinks (late) ----
\usepackage[colorlinks=true,linkcolor=accentdk,urlcolor=accent,
            citecolor=accentdk,linktoc=all]{hyperref}
\usepackage{url}
\urlstyle{sf}

%------------------------------------------------ markdown-style callouts ------
\usepackage[most]{tcolorbox}
\usetikzlibrary{shapes.geometric}   % 4-point "sparkle" star on the title page
\usepackage{graphicx}
\graphicspath{{./}{doc/}}
% Lime-tile icon. Uses the real DFFRNT PNG logo when `dffrnt-icon.png` is present
% beside this file; otherwise falls back to a vector sparkle so the page always
% renders. Drop the PNG in to get the true logo; delete this macro's fallback
% branch (leaving just the circle) to have empty tiles instead.
\newcommand{\dffrntlogo}[1]{%
  \IfFileExists{dffrnt-icon.png}
    {\node at (#1) {\includegraphics[width=15mm]{dffrnt-icon.png}};}
    {\node[icontile] at (#1) {};
     \node[sparkle]  at (#1) {};
     \node[star, star points=4, star point ratio=0.34, fill=brandlime,
           minimum size=2.6mm, inner sep=0] at ([xshift=-3.9mm,yshift=-3.4mm]#1) {};}%
}

\newtcolorbox{gnote}{
  enhanced, breakable,
  colback=notebg, colframe=noteblue,
  boxrule=0pt, leftrule=3.5pt, sharp corners,
  left=12pt, right=12pt, top=8pt, bottom=8pt,
  before skip=10pt, after skip=12pt,
  fontupper=\small\color{ink},
  before upper={{\sffamily\bfseries\color{noteblue}Note}\par\smallskip},
}
\newtcolorbox{gtip}{
  enhanced, breakable,
  colback=tipbg, colframe=tipgreen,
  boxrule=0pt, leftrule=3.5pt, sharp corners,
  left=12pt, right=12pt, top=8pt, bottom=8pt,
  before skip=10pt, after skip=12pt,
  fontupper=\small\color{ink},
  before upper={{\sffamily\bfseries\color{tipgreen}Tip}\par\smallskip},
}
\newtcolorbox{gwarn}{
  enhanced, breakable,
  colback=warnbg, colframe=warnred,
  boxrule=0pt, leftrule=3.5pt, sharp corners,
  left=12pt, right=12pt, top=8pt, bottom=8pt,
  before skip=10pt, after skip=12pt,
  fontupper=\small\color{ink},
  before upper={{\sffamily\bfseries\color{warnred}Warning}\par\smallskip},
}

\setcounter{tocdepth}{2}
\setcounter{secnumdepth}{3}

% Start each major section a little lower.
\let\oldsection\section
\renewcommand{\section}{\vspace{1em}\oldsection}

%==============================================================================
\begin{document}
%==============================================================================

%--------------------------------------------------------------- TITLE PAGE ----
\begin{titlepage}
\thispagestyle{empty}
\begin{tikzpicture}[remember picture, overlay,
    icontile/.style={fill=icnk, rounded corners=2mm, minimum size=15mm, inner sep=0},
    sparkle/.style={star, star points=4, star point ratio=0.36, fill=brandlime,
                    minimum size=8mm, inner sep=0}]
  % page background
  \fill[pagebg] (current page.north west) rectangle (current page.south east);

  % --- decorative circle cluster (top) ---
  \fill[accent]    ([xshift=83mm, yshift=-40mm]current page.north west) circle (17mm);
  \fill[accent]    ([xshift=167mm,yshift=-43mm]current page.north west) circle (17mm);
  \fill[accent]    ([xshift=131mm,yshift=-85mm]current page.north west) circle (15mm);
  \fill[accent]    ([xshift=170mm,yshift=-122mm]current page.north west) circle (14mm);
  \coordinate (limeA) at ([xshift=128mm,yshift=-38mm]current page.north west);
  \coordinate (limeB) at ([xshift=172mm,yshift=-86mm]current page.north west);
  \fill[brandlime] (limeA) circle (17mm);
  \fill[brandlime] (limeB) circle (17mm);
  \dffrntlogo{limeA}
  \dffrntlogo{limeB}

  % --- product title + tagline ---
  \node[anchor=west, text=accentdk,
        font=\fontfamily{phv}\bfseries\fontsize{40}{46}\selectfont]
    at ([xshift=26mm,yshift=-152mm]current page.north west) {DFFRNT};
  \node[anchor=west, text=accentdk,
        font=\fontfamily{phv}\bfseries\fontsize{40}{46}\selectfont]
    at ([xshift=26mm,yshift=-168mm]current page.north west) {AI Assistant};
  \node[anchor=north west, text=accent, font=\fontfamily{phv}\selectfont\large]
    at ([xshift=26mm,yshift=-180mm]current page.north west)
    {\parbox{155mm}{A fully air-gapped, private knowledge and content AI assistant}};

  % --- guide block ---
  \fill[accent, rounded corners=5mm]
    ([xshift=20mm,yshift=-190mm]current page.north west)
    rectangle ([xshift=196mm,yshift=-258mm]current page.north west);
  \node[anchor=west, text=brandlime,
        font=\fontfamily{phv}\bfseries\fontsize{42}{46}\selectfont]
    at ([xshift=32mm,yshift=-212mm]current page.north west) {INSTALLER GUIDE};
  \node[anchor=north west, text=blocklt, font=\fontfamily{phv}\selectfont\large]
    at ([xshift=32mm,yshift=-222mm]current page.north west)
    {\parbox{152mm}{Deploying the assistant on a target machine}};
  \node[anchor=north west, text=blocklt2, font=\fontfamily{phv}\selectfont\footnotesize]
    at ([xshift=32mm,yshift=-234mm]current page.north west)
    {\parbox{152mm}{A private, fully offline RAG assistant that ingests company
     documents and answers questions with grounded, citation-backed responses.
     Everything runs locally --- no data leaves the machine.}};
  \node[anchor=west, text=brandlime, font=\fontfamily{phv}\bfseries\selectfont\small]
    at ([xshift=32mm,yshift=-252mm]current page.north west) {\today};
\end{tikzpicture}
\end{titlepage}

%---------------------------------------------------------- FRONT MATTER -------
\pagenumbering{roman}
\setcounter{page}{1}
\pagestyle{fancy}

{\hypersetup{linkcolor=ink}\tableofcontents}
\clearpage

%------------------------------------------------------------- MAIN MATTER ------
\pagenumbering{arabic}
\setcounter{page}{1}

This guide walks you through installing the DFFRNT Air-Gapped AI Assistant on
a target machine, from receiving the installer bundle to a running system. It
is written for a general audience --- no prior experience with Docker or the
command line is assumed beyond being able to open a terminal and type
commands.

\section{Concepts and technologies}

You do not need to master any of these to install the assistant, but knowing
what each piece does makes the process (and any error messages) much less
mysterious.

\subsection{The assistant itself}

The DFFRNT AI Assistant is a private, fully offline ``RAG''
(Retrieval-Augmented Generation) system. It ingests company documents
(proposals, resumes, policies, spreadsheets, \ldots) and answers questions
about them with citation-backed responses through a web page in your browser.

\begin{gnote}
Everything runs locally: the language model, the document index, and the
application. \textbf{No data ever leaves the machine.} That is the whole point
of the air-gapped design --- the assistant makes no external calls once it is
running.
\end{gnote}

\subsection{Docker}

\href{https://www.docker.com/}{Docker} is software that runs applications
inside \textbf{containers} --- self-contained boxes that carry everything an
application needs (its code, libraries, and settings). Because a container
ships with everything included, the assistant behaves identically on any
machine, and the target machine needs almost nothing pre-installed ---
\textbf{only Docker itself}. No Python, no packages, no scripts.

Related terms you will encounter:

\begin{itemize}
  \item \textbf{Image} --- a frozen snapshot of an application, like an
        installation CD. The bundle ships images; installing ``loads'' them
        into Docker.
  \item \textbf{Container} --- a running copy of an image. Stopping the stack
        stops the containers; the images stay behind for the next start.
  \item \textbf{Docker Compose} --- a Docker feature that starts and connects
        several containers together as one ``stack''. The assistant's stack
        has three services: Qdrant, Ollama, and the API/UI.
\end{itemize}

\subsection{Ollama}

\href{https://ollama.com/}{Ollama} is a local server for large language
models (LLMs). It hosts two models for the assistant: the chat model (the
\texttt{qwen3} family) that writes the answers, and an embedding model
(\texttt{bge-m3}) that turns text into vectors for searching. Ollama runs in
its own container and is not visible to you day-to-day.

\subsection{Qdrant}

\href{https://qdrant.tech/}{Qdrant} is a \textbf{vector database}. When
documents are ingested, they are split into chunks and each chunk is stored
as a vector (a list of numbers capturing its meaning). When you ask a
question, Qdrant finds the chunks whose meaning is closest to the question,
and those chunks become the evidence the model answers from. Qdrant also runs
in its own container.

\subsection{dffrnt-manager}

\texttt{dffrnt-manager} is the single tool you interact with. It is one
self-contained program with two faces:

\begin{itemize}
  \item Run it \textbf{with no arguments} and it opens a control panel in
        your web browser (Install / Manage / Models / Logs tabs). This is the
        recommended path for most users.
  \item Run it \textbf{with a command} (\texttt{install}, \texttt{start},
        \texttt{status}, \ldots) and it does the same job headlessly in the
        terminal --- useful over SSH or on servers.
\end{itemize}

\subsection{OFFLINE vs.\ online (AWS) bundles}

The bundle comes in two flavours, chosen when it was built:

\begin{itemize}
  \item \textbf{OFFLINE} --- carries every container image \emph{and} the AI
        models inside the bundle. Installing needs \textbf{no internet at
        all}. This is the flavour for air-gapped machines. The bundle is
        large (several GB).
  \item \textbf{AWS / online} --- ships only the application image and
        downloads the rest (Qdrant, Ollama, models) from the internet during
        installation. Smaller bundle, but the target must be online for the
        first install.
\end{itemize}

The README inside the bundle states which flavour you have.

\section{Step-by-step installation}

\subsection{What you should end up with}

The installer is distributed as a single \texttt{.tar} archive. Unpacking it
yields exactly \textbf{three files, which must stay together in the same
folder}:

\begin{lstlisting}
dffrnt-offline.tar.gz    (or dffrnt-aws.tar.gz -- the application bundle)
dffrnt-manager           (the installer + control panel, one program)
README.md                (target-side notes generated at build time)
\end{lstlisting}

\subsection{Step 0 --- Check the prerequisites}

The target machine needs:

\begin{enumerate}
  \item \textbf{Docker Engine with the Compose v2 plugin.}
    \begin{itemize}
      \item Linux: install \texttt{docker} and
            \texttt{docker-compose-plugin} from your distribution or from
            \url{https://docs.docker.com/engine/install/}.
      \item Windows / macOS: install \textbf{Docker Desktop} and make sure it
            is running (whale icon in the system tray / menu bar).
      \item Verify with: \texttt{docker --version} and
            \texttt{docker compose version}.
    \end{itemize}
  \item \textbf{Disk space} --- allow at least 25\,GB free for an OFFLINE
        install (images, models, and your document index all live on disk).
  \item On Linux, your user should be able to run Docker without
        \texttt{sudo} (see Section~\ref{sec:troubleshooting} if
        \texttt{docker ps} says ``permission denied'').
\end{enumerate}

\begin{gtip}
Don't worry about getting this perfect --- the installer checks all
prerequisites first and reports exactly what is missing if a check fails,
\emph{without changing anything on the machine}. It is always safe to run and
re-run.
\end{gtip}

\subsection{Step 1 --- Get the archive onto the machine}

\textbf{Option A --- download from a GitHub release (online machine):}

\begin{lstlisting}
curl -L -o dffrnt-installer.tar \
  https://github.com/ashwathram/dffrnt-airgapped-ai-assistant/releases/download/<TAG>/dffrnt-installer.tar
\end{lstlisting}

Replace \texttt{<TAG>} with the release version you were given (e.g.\
\texttt{v1.0.0}). You can also simply download it from the release page in a
browser.

\textbf{Option B --- USB thumb drive (air-gapped machine):}

\begin{enumerate}
  \item On a networked machine, download the \texttt{.tar} as above and copy
        it to the drive.
  \item Plug the drive into the target machine and copy the \texttt{.tar} to
        somewhere like your home folder or Desktop.
\end{enumerate}

\begin{gwarn}
The USB drive must \textbf{not} be FAT32-formatted. FAT32 cannot hold files
larger than 4\,GB, and an OFFLINE bundle is bigger than that --- the copy will
fail partway. Use exFAT or NTFS (most drives sold as ``64\,GB or larger'' are
already exFAT).
\end{gwarn}

\subsection{Step 2 --- Unpack the archive}

In a terminal, in the folder holding the \texttt{.tar}:

\begin{lstlisting}
mkdir dffrnt-installer
tar -xf dffrnt-installer.tar -C dffrnt-installer
cd dffrnt-installer
ls
\end{lstlisting}

\texttt{ls} should show the three files listed above. If your archive
unpacked into a nested subfolder, \texttt{cd} into it --- the point is to run
the next command from the folder that directly contains
\texttt{dffrnt-manager} and the \texttt{.tar.gz}.

On Windows or macOS you can instead right-click the archive and choose
``Extract'' --- then open a terminal in the extracted folder.

\subsection{Step 3 --- Run the installer}

\textbf{Browser (recommended):}

\begin{lstlisting}
./dffrnt-manager
\end{lstlisting}

This opens the control panel in your browser (it serves only to this machine
--- nothing is exposed to the network). Go to the \textbf{Install} tab: it
auto-detects the \texttt{.tar.gz} sitting next to it; you pick the
destination folder and click Install.

\textbf{Terminal (headless / SSH):}

\begin{lstlisting}
./dffrnt-manager install                     # installs to ./dffrnt
./dffrnt-manager install --dest /opt/dffrnt  # or choose the destination
\end{lstlisting}

Either way, the installer:

\begin{enumerate}
  \item checks the prerequisites (and stops with a clear message if one
        fails),
  \item unpacks the application bundle into the destination folder,
  \item loads the container image(s) into Docker,
  \item starts the stack --- Qdrant, Ollama, and the API/UI.
\end{enumerate}

An OFFLINE install runs entirely from the bundle. An online (AWS) install
pulls the base images and models from the internet at this point, so the
first install can take a while depending on your connection.

\subsection{Step 4 --- Verify}

\begin{lstlisting}
cd dffrnt                    # or the destination you chose
./dffrnt-manager status
\end{lstlisting}

\texttt{status} reports the health of each service. When everything is up,
open \url{http://localhost:8000} in a browser --- that is the assistant.

\begin{gnote}
On every start the assistant pre-loads the models into memory in the
background. The very first question after a start may take a little longer
while that warmup finishes --- and on a fresh \emph{online} install, the first
warmup also waits for the model download, which can take a few minutes.
\end{gnote}

\subsection{Step 5 --- Day-to-day management}
\label{sec:daytoday}

A copy of \texttt{dffrnt-manager} was placed in the install folder; it
manages the stack from there. Running it with no arguments opens the panel
(Manage / Models / Logs tabs); the same verbs work in the terminal:

\begin{lstlisting}
cd dffrnt
./dffrnt-manager             # browser control panel
./dffrnt-manager start       # bring the stack up
./dffrnt-manager stop        # stop all containers
./dffrnt-manager restart
./dffrnt-manager status      # service + API health
./dffrnt-manager logs        # follow container logs (or: logs api)
./dffrnt-manager audit       # follow the app's audit trail
./dffrnt-manager reingest    # force re-ingest of every stored document
./dffrnt-manager models      # list / switch / import / export models
\end{lstlisting}

\section{Accessing the control panel over SSH (AWS)}

\textbf{You only need this section for a remote, headless box you administer
over SSH --- in practice, an AWS instance.} On any local machine --- a Linux,
Windows, or macOS desktop sitting in front of you --- running
\texttt{./dffrnt-manager} opens the panel in your browser directly, and none
of this applies.

Even on AWS it is \textbf{optional}: every management action is also a
headless CLI command (see Step~5, \S\ref{sec:daytoday}), so you can run the
whole deployment over a plain SSH session without ever opening the panel. Use
the tunnel below only if you prefer the graphical Install / Manage / Models /
Logs views.

\subsection{Why a tunnel is needed}

For safety the panel binds to \textbf{127.0.0.1 only} (the machine's own
loopback) and puts a random per-session token in the URL it prints.
Loopback-only means it is deliberately \emph{not} reachable across the network
--- so on a remote box you reach it by forwarding your local traffic into that
loopback over the SSH connection you already have.

\begin{gwarn}
Do \textbf{not} open port 8901 (or 8000) in the AWS security group for this.
The tunnel rides entirely inside SSH (port 22); the panel stays invisible to
the internet. Opening those ports would expose the unauthenticated management
surface and the app --- exactly what the loopback binding prevents.
\end{gwarn}

\subsection{Using a SOCKS proxy (\texttt{ssh -D})}

A dynamic SOCKS proxy is the tidiest option, because it lets your browser use
the \emph{exact} \texttt{http://127.0.0.1:PORT/\ldots} URL the panel prints ---
unchanged, token and all. The browser hands \texttt{127.0.0.1} to the proxy,
and SSH resolves it on the \emph{remote} loopback, which is where the panel
lives.

\textbf{1. On the AWS box}, start the panel without opening a browser, on a
fixed port so it's predictable:

\begin{lstlisting}
./dffrnt-manager panel --no-browser --port 8901
\end{lstlisting}

Copy the full URL it prints --- it includes the token, e.g.\
\texttt{http://127.0.0.1:8901/?token=ab12cd34\ldots}

\textbf{2. On your own machine}, open a SOCKS proxy over SSH and leave it
running:

\begin{lstlisting}
ssh -D 1080 -N user@your-aws-host
\end{lstlisting}

\texttt{-D 1080} opens a SOCKS proxy on local port 1080; \texttt{-N} says
``just forward, don't run a remote shell''. (Add \texttt{-f} to push it into
the background.)

\textbf{3. Point your browser at the proxy:}

\begin{itemize}
  \item \textbf{Firefox:} Settings $\rightarrow$ Network Settings $\rightarrow$
        \emph{Manual proxy configuration} $\rightarrow$ SOCKS Host
        \texttt{127.0.0.1}, Port \texttt{1080}, select \textbf{SOCKS v5}, and
        tick \textbf{Proxy DNS when using SOCKS v5}.
  \item \textbf{Chrome / Chromium:} launch it with
        \texttt{-{}-proxy-server="socks5://127.0.0.1:1080"}, or use a
        proxy-switcher extension so you don't reroute your everyday browsing.
\end{itemize}

\textbf{4. Open the panel.} Paste the exact URL the box printed
(\texttt{http://127.0.0.1:8901/?token=\ldots}) into that browser. The request
tunnels to AWS, where \texttt{127.0.0.1:8901} is the panel --- you now have the
full control panel as if you were sitting at the machine.

When you're done, remove the browser's proxy setting and stop the SSH session
(Ctrl-C, or \texttt{kill} it if you used \texttt{-f}).

\begin{gnote}
The SOCKS proxy reroutes \emph{all} of that browser's traffic through the AWS
box while it's set. Use a separate browser profile or a proxy-switcher
extension, and remember to turn the proxy off afterwards --- otherwise normal
browsing will stall once the tunnel closes.
\end{gnote}

\subsection{Alternative: local port forwarding (\texttt{ssh -L})}

If you'd rather not touch any browser proxy setting, forward a single local
port to the panel instead:

\begin{lstlisting}
ssh -L 8901:127.0.0.1:8901 -N user@your-aws-host
\end{lstlisting}

Then open \texttt{http://127.0.0.1:8901/?token=\ldots} (the same printed URL)
in your normal browser --- no proxy configuration needed. The catch is that
the local port must match the panel's port, which is why we fixed it with
\texttt{-{}-port 8901} above.

\section{FAQ}

\paragraph{Do I need to know Docker to use this?}
No. \texttt{dffrnt-manager} drives Docker for you --- install, start, stop,
status, and logs are all one command or one click. Docker just needs to be
installed and running.

\paragraph{How do I see what's running?}
\texttt{./dffrnt-manager status} is the friendly view. If you're curious
about the underlying containers, \texttt{docker ps} lists them (you'll see
three: qdrant, ollama, and the api), and Docker Desktop shows the same thing
graphically on Windows/macOS. You should not need to touch containers
directly --- always prefer the manager, which starts and stops them in the
right order.

\paragraph{Does the machine need internet?}
For an \textbf{OFFLINE} bundle: no, not even during install. For an
\textbf{AWS/online} bundle: yes, once, during install (and when pulling a new
model later). After that, the assistant itself never calls out --- everything
runs locally.

\paragraph{Where do my documents and data live?}
Inside the install folder: the vector index in \texttt{qdrant\_storage/}, the
AI models in \texttt{ollama\_models/}, uploaded documents and app data in
\texttt{data/}, and logs in \texttt{logs/}. Deleting the install folder
deletes your data.

\paragraph{How do I stop it / start it after a reboot?}
\texttt{./dffrnt-manager stop} and \texttt{./dffrnt-manager start} from the
install folder. The stack does not start automatically on boot unless you set
that up yourself.

\paragraph{What is the address the panel opens, and why does it look odd?}
The control panel serves on \texttt{127.0.0.1} (this machine only) with a
one-time token in the URL --- that's a safety measure so nothing on the
network can reach the management functions. The assistant itself is the
friendlier \url{http://localhost:8000}.

\paragraph{Can other people on the network use the assistant?}
If they can reach port 8000 on this machine, yes.

\begin{gwarn}
The assistant has \textbf{no login}. Anyone who can reach port 8000 can query,
upload, and delete documents. Only expose it to networks you trust --- on a
server, restrict the firewall / security group for port 8000 to your VPN or
office network, or put an authenticating reverse proxy in front. Never open it
to the whole internet.
\end{gwarn}

\paragraph{How do I change the AI model?}
Panel $\rightarrow$ \textbf{Models} tab, or \texttt{./dffrnt-manager models
use qwen3:8b}. On an air-gapped machine, first bring the model over as a
tarball: \texttt{models export} on any networked machine with this app, then
\texttt{models import} here. Both are in the Models tab too.

\paragraph{Can I change settings?}
Yes --- \texttt{config.toml} in the install folder holds every setting
(models, ports, GPU on/off, retrieval tuning). Edit it, then
\texttt{./dffrnt-manager restart}.

\paragraph{Will reinstalling wipe my settings?}
No. An existing edited \texttt{config.toml} is kept by default; the installer
only replaces it if you explicitly pass \texttt{--overwrite-config} (the old
file is then archived as \texttt{config.toml.old}).

\paragraph{Does anything get sent to the cloud?}
No. The models run locally under Ollama, the index is local Qdrant, and the
app makes no external calls. That is the point of the air-gapped design.

\section{Troubleshooting}
\label{sec:troubleshooting}

\subsection{``permission denied'' when the installer (or \texttt{docker ps}) talks to Docker}

On Linux, your user isn't in the \texttt{docker} group. Fix:

\begin{lstlisting}
sudo usermod -aG docker $USER
\end{lstlisting}

then \textbf{log out and back in} (or run \texttt{newgrp docker} in the
current terminal). Verify with \texttt{docker ps} --- it should list
containers (or print an empty table) without \texttt{sudo}.

\subsection{\texttt{./dffrnt-manager: Permission denied}}

The executable bit was lost in transfer (common when the file passed through
a FAT32 drive or some download tools). Fix:

\begin{lstlisting}
chmod +x dffrnt-manager
\end{lstlisting}

\subsection{\texttt{cannot execute binary file} or \texttt{Exec format error}}

The bundle was built for a different OS or CPU architecture than this machine
(e.g.\ an x86\_64 build on an ARM machine). Bundles are platform-specific ---
ask for a build matching the target and check with \texttt{uname -m}.

\subsection{macOS blocks the program (``cannot be opened because the developer cannot be verified'')}

macOS quarantines files that arrived from the internet. Clear it with:

\begin{lstlisting}
xattr -d com.apple.quarantine dffrnt-manager
\end{lstlisting}

or right-click $\rightarrow$ Open the first time.

\subsection{``No bundle found'' / installer can't find the tarball}

The three files must sit \textbf{together in the same folder}, and you must
run \texttt{./dffrnt-manager} from that folder. Check with \texttt{ls} that
you see the \texttt{dffrnt-*.tar.gz} next to \texttt{dffrnt-manager} --- if
extraction created a nested subfolder, \texttt{cd} into it. You can also name
the bundle explicitly:
\texttt{./dffrnt-manager install path/to/dffrnt-offline.tar.gz}.

\subsection{A prerequisite check fails}

The message names exactly what's missing --- most commonly Docker isn't
installed, isn't running (start Docker Desktop, or \texttt{sudo systemctl
start docker} on Linux), or the Compose v2 plugin is absent (\texttt{docker
compose version} should print a version). Install/start the missing piece and
re-run; the check changed nothing, so it's safe to retry.

\subsection{Port 8000 is already in use}

Another program on the machine owns port 8000. Edit \texttt{api\_port} in
\texttt{config.toml} in the install folder (e.g.\ \texttt{api\_port = 8080}),
then \texttt{./dffrnt-manager restart}, and use
\url{http://localhost:8080}.

\subsection{\texttt{status} says a service is unhealthy / the UI won't load}

Give it a minute after \texttt{start} --- the models take time to load. Then:

\begin{lstlisting}
./dffrnt-manager status      # which service is degraded?
./dffrnt-manager logs api    # read that service's log (also: logs ollama, logs qdrant)
./dffrnt-manager restart     # often clears one-off startup races
\end{lstlisting}

\subsection{GPU mode refuses to start}

\texttt{gpu = true} in \texttt{config.toml} requires the NVIDIA driver
\textbf{and} the NVIDIA Container Toolkit on the host; the manager checks and
refuses with a message if either is missing. Install them (e.g.\
\texttt{nvidia-container-toolkit} from your distro or NVIDIA's repository),
or set \texttt{gpu = false} to run on CPU.

\subsection{The USB drive won't take the bundle / a model export}

The drive is FAT32, whose 4\,GB per-file limit is smaller than the bundle and
most models. Reformat the drive as exFAT (both Windows and macOS/Linux can
read and write it) and copy again.

\subsection{Disk fills up over time}

Model switching is designed not to accumulate: switching removes the
previously configured model from the cache. If space runs low anyway, check
\texttt{./dffrnt-manager models} for unused models (\texttt{models rm NAME})
and \texttt{docker system df} for Docker's own usage.

\subsection{Something else}

\texttt{./dffrnt-manager logs} (all services) and
\texttt{logs/installer.log} in the install folder record what happened ---
include the relevant lines when asking for help.

\end{document}
